\documentclass[11pt]{article}
\usepackage[margin=1in]{geometry}
\usepackage{amsmath,amssymb,bm}
\newcommand{\bten}{\bm{b}}
\newcommand{\xiv}{\bm{\xi}}
\newcommand{\Ften}{\bm{F}}
\usepackage{graphicx}
\usepackage{booktabs}
\usepackage{xcolor}
\usepackage[hidelinks]{hyperref}
\usepackage{listings}
\usepackage{enumitem}

\definecolor{kw}{rgb}{0.10,0.30,0.65}
\definecolor{cm}{rgb}{0.30,0.50,0.30}
\definecolor{st}{rgb}{0.60,0.20,0.20}
\lstdefinestyle{py}{language=Python,basicstyle=\ttfamily\small,breaklines=true,
  frame=single,backgroundcolor=\color{black!4},columns=fullflexible,
  showstringspaces=false,keywordstyle=\color{kw},commentstyle=\color{cm},
  stringstyle=\color{st}}
\lstdefinestyle{sh}{basicstyle=\ttfamily\small,breaklines=true,frame=single,
  backgroundcolor=\color{black!4},columns=fullflexible}
\lstset{style=py}

\title{\textbf{pyRMT --- User Manual}\\[2mm]
\large Installing, running, and building simulations with the
Reference Map Technique solver}
\author{Saman Seifi}
\date{\today}

\begin{document}
\maketitle
\begin{abstract}
\noindent
This manual is the practical guide to pyRMT: how to install it, run the supplied
benchmarks, use the public API, build a custom 2D fluid--structure-interaction
simulation, and interpret the outputs. For the underlying method and equations see
the companion \emph{Technical Manual}.
\end{abstract}
\tableofcontents

\section{Overview}
pyRMT simulates 2D incompressible fluid--structure interaction (FSI) with the
Reference Map Technique: a hyperelastic or viscoelastic solid and a Newtonian fluid
share one velocity field on a fixed Eulerian grid. It provides two solver families:
\begin{itemize}[nosep]
  \item \textbf{Staggered / MAC} (\texttt{pyRMT/mac.py}) --- exact divergence-free
  projection; recommended for surface tension, contact, and viscoelasticity.
  \item \textbf{Collocated} (\texttt{pyRMT/functions.py}) --- Rhie--Chow; the
  original solver, with a variable-density (AMG) option.
\end{itemize}

\section{Installation}
\paragraph{Requirements.} Python $\ge 3.10$ and \texttt{numpy}, \texttt{scipy},
\texttt{matplotlib}, \texttt{h5py}, \texttt{pyamg}, \texttt{numba}. Optional:
\texttt{scikit-fmm} (fast-marching level-set reinitialization), \texttt{imageio}
(animated GIF benchmarks), \texttt{pytest} (tests).
\begin{lstlisting}[style=sh]
git clone https://github.com/samanseifi/pyRMT.git
cd pyRMT
pip install -e ".[test]"     # editable install + pytest + scikit-fmm
# minimal:           pip install -e .
# with FMM only:     pip install -e ".[fmm]"
\end{lstlisting}
The Numba-JIT kernels compile on first use (a one-time startup cost; caching is
enabled). The branching model is: develop from \texttt{develop}; \texttt{main} is the
stable branch.

\section{Repository layout}
\begin{center}
\begin{tabular}{ll}
\toprule
path & contents \\
\midrule
\texttt{pyRMT/} & the importable package (see Technical Manual, Module map) \\
\texttt{benchmarks/} & runnable cases (validation, demos, animations) \\
\texttt{tests/} & \texttt{pytest} suite \\
\texttt{docs/} & this manual, the technical manual, figures (\texttt{docs/img/}) \\
\texttt{outputs/} & generated CSV / figures / GIFs / frame data \\
\bottomrule
\end{tabular}
\end{center}

\section{Quick start}
Run any benchmark as a script; outputs land under \texttt{outputs/}.
\begin{lstlisting}[style=sh]
# --- pure fluid validation (collocated) ---
python benchmarks/lid_driven_cavity.py 1000           # Re=1000 vs Ghia

# --- MAC solver ---
python benchmarks/mac_lid_driven.py                   # lid cavity (staggered)
python benchmarks/mac_taylor_green_convergence.py     # 2nd-order fluid
python benchmarks/mac_disc_in_tg_convergence.py       # FSI convergence study
python benchmarks/mac_soft_disc_lid.py                # soft disc vs Sugiyama
python benchmarks/mac_surface_tension_relax.py        # square -> circle

# --- contact & multi-body (GIFs; need imageio) ---
python benchmarks/mac_two_disc_contact.py
python benchmarks/mac_multi_shape_showcase.py 256 12.0      # 6 shapes, two views
python benchmarks/mac_settling_shapes_gif.py

# --- viscoelasticity ---
python benchmarks/viscoelastic_relaxation_test.py          # 0-D verification
python benchmarks/mac_viscoelastic_uniform.py 96 6.0 0.5   # extension, tau=0.5
python benchmarks/mac_viscoelastic_uniform.py plot         # money figure
\end{lstlisting}

\section{Benchmark catalogue}
Most MAC drivers take command-line arguments \texttt{[N] [t\_end] [\dots]} and print
their resolution / parameters on startup. Selected cases:
\begin{center}
\small
\begin{tabular}{lll}
\toprule
script & what it shows & typical call \\
\midrule
\texttt{mac\_lid\_driven.py} & lid cavity (Ghia) & \texttt{... [Re]} \\
\texttt{mac\_taylor\_green\_convergence.py} & 2nd-order fluid & \texttt{...} \\
\texttt{mac\_disc\_in\_tg\_convergence.py} & FSI order, all schemes & \texttt{...} \\
\texttt{mac\_soft\_disc\_lid.py} & soft disc vs Sugiyama & \texttt{... [N] [scheme] [t]} \\
\texttt{mac\_surface\_tension\_drop.py} & static Laplace jump & \texttt{... [N]} \\
\texttt{mac\_surface\_tension\_relax.py} & square $\to$ circle & \texttt{... [N]} \\
\texttt{mac\_two\_disc\_contact.py} & collision rebound vs $\eta$ & \texttt{...} \\
\texttt{mac\_multi\_shape\_showcase.py} & 6 shapes, speed/stress/contact & \texttt{... N t [mode]} \\
\texttt{mac\_settling\_shapes\_gif.py} & gravity settling pile & \texttt{... [N] [t]} \\
\texttt{mac\_reinit\_test.py} & reference-map reinit & \texttt{... [N] [t] [reinit]} \\
\texttt{mac\_viscoelastic\_uniform.py} & extension, exact plateau & \texttt{... N t tau} \\
\texttt{viscoelastic\_relaxation\_test.py} & 0-D constitutive checks & \texttt{...} \\
\bottomrule
\end{tabular}
\end{center}
\texttt{tau<=0} selects the elastic ($\tau\!\to\!\infty$) limit; \texttt{scheme} is
one of \texttt{semilagrangian}, \texttt{semilagrangian\_cubic}, \texttt{central2},
\texttt{weno5}, \texttt{conservative}. The showcase \texttt{mode} is
\texttt{speed}/\texttt{stress}/\texttt{contact}; passing \texttt{rerender N} re-renders
saved frames without re-simulating.

\section{Public API tour}
After \texttt{pip install -e .} the package is importable as \texttt{pyRMT}.
\subsection{Grid, operators, projection (MAC)}
\begin{lstlisting}
import numpy as np
from pyRMT.mac import (mac_grid, poisson_eigs_neumann, project,
                       momentum_predictor, divergence)

N = 128
dx, dy = mac_grid(N, N)                 # spacing on the unit square
u = np.zeros((N, N+1))                  # x-velocity on x-faces
v = np.zeros((N+1, N))                  # y-velocity on y-faces
eig = poisson_eigs_neumann(N, N, dx, dy)  # precompute DCT eigenvalues

# one lid-driven predictor + exact projection step
ustar, vstar = momentum_predictor(u, v, nu=0.01, dx=dx, dy=dy, dt=1e-3,
                                   U_lid=1.0)
u, v, p = project(ustar, vstar, dx, dy, dt=1e-3, rho=1.0, eig=eig)
print("max|div| =", np.abs(divergence(u, v, dx, dy)).max())   # ~1e-14
\end{lstlisting}

\subsection{Reference map, stress, advection}
\begin{lstlisting}
from pyRMT.functions import (extrapolate_reference_map, advect_reference_map,
    rebuild_phi_from_reference_map, solid_cauchy_stress, smoothed_heaviside)

# initialise a disc: phi<0 inside
xc = (np.arange(N)+0.5)*dx; X, Y = np.meshgrid(xc, xc)
pin = lambda Xa, Ya: np.sqrt((Xa-0.5)**2 + (Ya-0.5)**2) - 0.2
phi = pin(X, Y); m = (phi <= 0).astype(float)
X1, X2 = extrapolate_reference_map(X*m, Y*m, phi, dx, dy, 3)   # reference map

# advect one step (cell-centred velocity u_c, v_c), then re-extrapolate
X1 = advect_reference_map(X1, u_c, v_c, Xg, Yg, dt, dx, dy, phi, 'semilagrangian', 0.0)*m
X1, X2 = extrapolate_reference_map(X1, X2, phi, dx, dy, 3)
phi = rebuild_phi_from_reference_map(X1, X2, pin)

# neo-Hookean solid stress, blended into the momentum force
sxx, sxy, syy, J = solid_cauchy_stress(X1, X2, dx, dy, mu_s=1.0, kappa=0.0, phi=phi)
H = smoothed_heaviside(phi, 2*dx)
Sxx = (1-H)*sxx     # (1-H) is the solid indicator
\end{lstlisting}

\subsection{Viscoelastic stress (log-conformation)}
\begin{lstlisting}
from pyRMT.viscoelastic import logconf_local_step, sym_exp, sym_log

# psi = log(b_e); start unstressed
p11 = np.zeros((N,N)); p12 = np.zeros((N,N)); p22 = np.zeros((N,N))
# advect psi like the reference map, then relax + stretch locally:
p11, p12, p22 = logconf_local_step(p11, p12, p22,
                                   L11, L12, L21, L22, tau=0.5, dt=dt)
be11, be12, be22 = sym_exp(p11, p12, p22)        # SPD by construction
Sxx = (1-H) * G * be11                            # sigma = G b_e
\end{lstlisting}

\subsection{Contact and surface tension}
\begin{lstlisting}
from pyRMT.mac import contact_stress, interfacial_force_faces
# trace-free contact stress between two solids (add to the stress tensor)
txx, txy, tyy = contact_stress(phi_a, phi_b, eta=2.5, Gsum=2*mu_s,
                               eps=3*dx, dx=dx, dy=dy)
# balanced-force surface tension as face forces
fu, fv = interfacial_force_faces(kappa, H, gamma=1.0, dx=dx, dy=dy)
\end{lstlisting}

\section{Building a custom simulation (MAC)}
A minimal soft-disc-in-flow loop:
\begin{lstlisting}
import numpy as np
from pyRMT.mac import (mac_grid, poisson_eigs_neumann, project, momentum_predictor)
from pyRMT.functions import (extrapolate_reference_map, advect_reference_map,
    rebuild_phi_from_reference_map, solid_cauchy_stress, smoothed_heaviside,
    grad_central_x_2nd, grad_central_y_2nd)

N, mu_s, nu, rho, U = 128, 1.0, 0.01, 1.0, 1.0
dx, dy = mac_grid(N, N)
xc = (np.arange(N)+0.5)*dx; X, Y = np.meshgrid(xc, xc)
Xg, Yg = np.meshgrid(np.arange(N)*dx, np.arange(N)*dy)
eig = poisson_eigs_neumann(N, N, dx, dy)
u = np.zeros((N, N+1)); v = np.zeros((N+1, N))
pin = lambda Xa, Ya: np.sqrt((Xa-0.5)**2+(Ya-0.5)**2) - 0.2
phi = pin(X, Y); m = (phi <= 0).astype(float)
X1, X2 = extrapolate_reference_map(X*m, Y*m, phi, dx, dy, 3)
dt = min(0.3*dx/U, 0.2*dx*dx/nu, 0.3*dx/np.sqrt(mu_s/rho))

for step in range(2000):
    u_c = 0.5*(u[:, :-1]+u[:, 1:]); v_c = 0.5*(v[:-1, :]+v[1:, :])
    phi = rebuild_phi_from_reference_map(X1, X2, pin); m = (phi <= 0).astype(float)
    X1 = advect_reference_map(X1, u_c, v_c, Xg, Yg, dt, dx, dy, phi, 'semilagrangian', 0.)*m
    X2 = advect_reference_map(X2, u_c, v_c, Xg, Yg, dt, dx, dy, phi, 'semilagrangian', 0.)*m
    X1, X2 = extrapolate_reference_map(X1, X2, phi, dx, dy, 3)

    sxx, sxy, syy, J = solid_cauchy_stress(X1, X2, dx, dy, mu_s, 0.0, phi)
    H = smoothed_heaviside(phi, 2*dx)
    Sxx, Sxy, Syy = (1-H)*sxx, (1-H)*sxy, (1-H)*syy
    divx = grad_central_x_2nd(Sxx, dx) + grad_central_y_2nd(Sxy, dy)
    divy = grad_central_x_2nd(Sxy, dx) + grad_central_y_2nd(Syy, dy)
    fu = np.zeros((N, N+1)); fu[:, 1:-1] = 0.5*(divx[:, 1:]+divx[:, :-1])
    fv = np.zeros((N+1, N)); fv[1:-1, :] = 0.5*(divy[1:, :]+divy[:-1, :])

    ustar, vstar = momentum_predictor(u, v, nu, dx, dy, dt, U, fu=fu, fv=fv, rho=rho)
    u, v, p = project(ustar, vstar, dx, dy, dt, rho, eig)
\end{lstlisting}
To make it viscoelastic, carry $\bm\psi=\log\bten_e$ (three cell-centred arrays),
advect them like \texttt{X1/X2}, call \texttt{logconf\_local\_step} with the velocity
gradient, set the stress to \texttt{(1-H)*G*be} via \texttt{sym\_exp}. For multiple
solids, keep one $(\xiv,\phi)$ per body and add \texttt{contact\_stress} for every
pair plus the walls. For periodic or free-slip boxes use
\texttt{momentum\_predictor\_periodic}/\texttt{\_freeslip} with the matching
projection.

\section{Outputs}
Benchmarks write to \texttt{outputs/<case>/}: CSV time series, PNG figures, and
(for the demos) animated GIFs. Long animations \emph{stream} per-frame field data to
\texttt{outputs/<case>/frames/} so they can be re-rendered cheaply
(\texttt{... rerender N}) without re-simulating. The curated figures used in the
documentation live under \texttt{docs/img/}.

\section{Testing}
\begin{lstlisting}[style=sh]
pytest                 # runs tests/ (operators, projection, Poisson, stress,
                       # contact, reinit, interpolation, viscoelastic)
pytest tests/test_viscoelastic.py -q
\end{lstlisting}
The suite checks the discrete operators and exact projection, the neo-Hookean
stress, contact, reinitialization, and the viscoelastic constitutive update against
analytic responses.

\section{Performance and memory notes}
\begin{itemize}[nosep]
  \item Numba JIT: first call compiles (cached). The bilinear/bicubic interpolators
  are deliberately non-parallel to avoid a Numba parallel-codegen issue.
  \item Direct Poisson solves are $O(N\log N)$; the projection is the cheap part.
  \item Time step is CFL-limited by the elastic-wave speed $\sqrt{\mu_s/\rho}$ and
  viscosity; stiff/fast solids force small $\Delta t$.
  \item Memory: holding all frames of a high-resolution animation in RAM can exceed
  a small machine; the showcase driver streams frames to disk to keep memory flat.
\end{itemize}

\section{Troubleshooting}
\begin{center}
\small
\begin{tabular}{p{0.46\textwidth}p{0.46\textwidth}}
\toprule
symptom & remedy \\
\midrule
\texttt{ImportError: scikit-fmm} & \texttt{pip install scikit-fmm} (or avoid FMM reinit) \\
GIF benchmark errors on save & \texttt{pip install imageio} \\
run dies silently, no error & likely OOM on a small box --- lower $N$, stream frames \\
solid ``folds'' (\,$\min J<0$\,) under sustained shear & enable reference-map reinit,
or use the viscoelastic model; for contact crushing, increase $N$ \\
near-contacting bodies stick / pinch & resolution-limited penalty contact; refine or
soften the lid (lubrication model planned) \\
FSI fields converge below 2nd order & expected: the diffuse interface caps the order \\
slow first run & Numba compilation; subsequent runs use the cache \\
\bottomrule
\end{tabular}
\end{center}

\section{Building this documentation}
\begin{lstlisting}[style=sh]
cd docs
pdflatex user_manual.tex
pdflatex technical_manual.tex
\end{lstlisting}

\end{document}
